%% ma: L10-B12-0046
%% lop: 10
%% chuong: 5
%% bai: 12
%% dang: 10.12.1
%% loai: TN4
%% muc: NB
%% dap_an: "C"
%% nguon: "chuyen-de-cac-so-dac-trung-cua-mau-so-lieu-khong-ghep-nhom-toan-10-chuong-trinh-moi.pdf (GV. Phan Nhật Linh) (P01), Bài 01 Phần I Câu 29"
%% kien_thuc: "Lớp 10 Bài 12 (số gần đúng, sai số tuyệt đối, độ chính xác, sai số tương đối, quy tròn số)"
%% trang_thai: chinh-thuc
%% ghi_chu: "Tự duyệt theo yêu cầu giáo viên 22/10/2026."
\begin{ex}
Độ cao của một ngọn núi được ghi lại là $\overline h=1372{,}5\text{ m}\pm0{,}2\text{ m}$. Độ chính xác $d$ của phép đo trên là
\choice{$d=0{,}1$ m}{$d=1$ m}{\True $d=0{,}2$ m}{$d=2$ m}
\loigiai{Theo cách viết $\overline h=a\pm d$, độ chính xác $d=0{,}2$ m. Chọn C.}
\end{ex}
